% Shared exercise chunk: l1-stable-softmax
\begin{exercisechunk}
\item \textbf{Stable softmax and cross-entropy from logits.}\suppmark
\label[exercise]{ex:l1-stable-softmax}
\exerciseprofile{3}{3}{3}{\exproof\quad\excalculation\quad\exinterpretation}{%
Derive logit-based softmax and cross-entropy calculations that preserve small
losses and gradients and avoid subtracting nearly equal probabilities.}
Let \(K\ge 2\), let \(z\in\R^K\), and define
\[
 q_k=\operatorname{softmax}(z)_k=\frac{e^{z_k}}{\sum_{i=1}^K e^{z_i}}.
\]
Write
\[
 \operatorname{LSE}(a_1,\ldots,a_m)
 =\log\!\left(\sum_{i=1}^m e^{a_i}\right),
 \qquad
 \operatorname{softplus}(a)=\log(1+e^a).
\]
Derive all identities first in exact real arithmetic. In floating-point
calculations, assume that \(\operatorname{LSE}\), \(\operatorname{softplus}\),
\(\operatorname{log1p}(x)=\log(1+x)\), and
\(\operatorname{expm1}(x)=e^x-1\) are evaluated by stable library routines.

\begin{enumerate}[(a),beginpenalty=10000]
\item Prove that
\(\operatorname{softmax}(z-c\mathbf 1)=\operatorname{softmax}(z)\) for every
\(c\in\R\). For \(c=\max_i z_i\), derive
\[
 q_k=\frac{e^{z_k-c}}{\sum_{i=1}^K e^{z_i-c}},
 \qquad
 \log q_k=(z_k-c)-\operatorname{LSE}(z_1-c,\ldots,z_K-c).
\]
Explain why the shift prevents overflow in the exponentials and why one
should compute log probabilities from the second formula instead of taking
the logarithm of a rounded probability.

\item Fix a hard target \(t\in[K]\), let \(e_t\in\R^K\) be the \(t\)-th
standard basis vector, and define
\[
 a=\log\!\left(\sum_{j\ne t}e^{z_j-z_t}\right)
 =\operatorname{LSE}\bigl((z_j-z_t)_{j\ne t}\bigr).
\]
Show that the hard-target cross-entropy and its logit gradient can be written
as
\[
 L(z,t)=-\log q_t=\operatorname{softplus}(a),
\]
and
\[
 g_j=\frac{\partial L}{\partial z_j}
 =e^{z_j-z_t-L(z,t)}\quad(j\ne t),
 \qquad
 g_t=-\sum_{j\ne t}g_j.
\]
Conclude that \(\sum_{k=1}^K g_k=0\). Explain how stable
\(\operatorname{LSE}\) and \(\operatorname{softplus}\) routines evaluate the
loss without forming a large exponential, and why the displayed formula for
\(g_t\) avoids subtracting a probability close to \(1\) from \(1\).

\item Consider \(z=(9,-9,-9)\) and \(t=1\) in float32 arithmetic. Put
\(r=e^{-18}\approx1.523\times10^{-8}\). Assume that \(r\) remains nonzero
but round-to-nearest gives
\[
 \operatorname{fl}(1+2r)=1.
\]
The usual maximum shift has already been applied. Compute the loss, gradient,
and sum of the gradient coordinates when an implementation forms \(q\) and
then uses \(L=-\log q_t\) and \(g=q-e_t\). Compute the corresponding
quantities from part~(b), to first order in \(r\). Identify the numerical
invariant that the usual calculation violates.

\item Now let \(z_{\rm S},z_{\rm T}\in\R^K\) be student and teacher logits,
and put
\[
 q=\operatorname{softmax}(z_{\rm S}),
 \qquad p=\operatorname{softmax}(z_{\rm T}),
 \qquad \delta=z_{\rm S}-z_{\rm T}.
\]
The exact soft-target cross-entropy gradient with respect to \(z_{\rm S}\)
is \(q-p\). First prove
\[
 \log\frac{q_k}{p_k}
 =\delta_k-\log\!\left(\sum_{i=1}^Kp_i e^{\delta_i}\right).
\]
\par\noindent
Define
\[
 \bar\delta=\sum_{i=1}^Kp_i\delta_i,
 \qquad d_k=\delta_k-\bar\delta,
\]
\[
 s_k=d_k-
 \operatorname{log1p}\!\left(
   \sum_{i=1}^Kp_i\operatorname{expm1}(d_i)
 \right),
 \qquad
 u_k=p_k\operatorname{expm1}(s_k).
\]
Show that
\(\sum_i p_i\operatorname{expm1}(d_i)\ge0\), so the displayed
\(\operatorname{log1p}\) has a nonnegative argument.
Prove that \(s_k=\log(q_k/p_k)\) and \(u_k=q_k-p_k\). Verify directly
that \(u=0\) when \(z_{\rm S}=z_{\rm T}\), and explain why this calculation
avoids subtracting two nearly equal rounded probabilities. State why an
implementation needs the teacher logits, rather than only the rounded teacher
probabilities, to use this calculation.

\item \emph{Optional extension.} A floating-point evaluation of \(u\) may
have a small nonzero sum. Define
\[
 \widehat g_k=u_k-q_k\sum_{i=1}^K u_i.
\]
Show that \(\sum_k\widehat g_k=0\) whenever \(\sum_kq_k=1\). Explain why
the correction leaves the exact real-arithmetic gradient unchanged.
\end{enumerate}

\begin{hints}
For (b), factor \(e^{z_t}\) from the softmax denominator and differentiate
through the differences \(z_j-z_t\). For (d), write the student normalizer in
terms of the teacher normalizer and use
\(\sum_i p_i e^{d_i}=1+\sum_i p_i\operatorname{expm1}(d_i)\).
\end{hints}
\end{exercisechunk}
